\section{Introduction}

A safety evaluation is only worth running if the model cannot tell it is one.
There is now direct evidence that this assumption is breaking down. In Anthropic's
reporting on Claude Sonnet~4.5, the model verbalized awareness that it was being
tested in roughly $16\%$ of code-safety evaluation transcripts and roughly $26\%$
of SWE-bench-style transcripts, against well under $1\%$ of real user traffic; on a
political-sycophancy probe the model explicitly flagged the test framing in about
$13\%$ of transcripts~\cite{anthropic2025sonnet45}. When a model behaves
differently because it has recognized the evaluation frame, the measured behavior
is no longer the deployed behavior, and the entire chain of safety arguments built
on those measurements is compromised.

This motivates a question that, to our knowledge, has not been posed in exactly
this form: \emph{given an internal ``I am being tested'' concept, can we prove
which inputs reach it and which provably cannot?} A proof of \emph{unreachability}
around an activation would be a genuine guarantee --- not a correlation, not an
empirical attack that simply failed to find anything --- that no small perturbation
can flip the model into the eval-aware state. A constructive \emph{reachability}
witness, conversely, would be a concrete prompt that demonstrably crosses into the
state. Both are useful: the first underwrites a claim that a particular evaluation
context is indistinguishable from deployment; the second is a red-team tool.

\paragraph{A two-paper mash-up.}
We obtain a working answer for a small system by combining two prior lines of work
that were not designed to meet.

\emph{Defining the concept (the NLA half).} Anthropic's Natural Language
Autoencoder~\cite{nla2026} verbalizes a model's internal activations into text. We
repurpose this idea to \emph{define} the target concept: rather than hand-label
activations, we let a verbalizer turn an activation into an English description and
let a text classifier decide whether that description reads as ``the model knows it
is being evaluated.'' The concept is thus grounded in language the model's own
internal state emits, not in an externally imposed direction.

\emph{Certifying reachability (the Raghunathan half).} Raghunathan, Steinhardt and
Liang~\cite{raghunathan2018certified,raghunathan2018sdp} introduced a sound convex
relaxation that produces a \emph{certificate} of robustness for a one-hidden-layer
ReLU network: an upper bound on the worst-case output over an $L_\infty$ ball that,
crucially, is itself differentiable and can be optimized. We use exactly this
machinery, but turned around: instead of certifying that a classifier's label
\emph{cannot} be flipped, we certify that the eval-aware region \emph{cannot} be
reached from a benign activation within a perturbation budget.

The bridge between the two halves is that we make the certified object a
one-hidden-layer ReLU network on continuous activations --- precisely Raghunathan's
architecture class. Since the relaxation cannot be pushed soundly through a
discrete text decoder and an n-gram classifier, we distil the NLA-text-to-classifier
map into such a network (a \emph{surrogate} $g$) and report how faithfully $g$
tracks the classifier rather than pretending the distillation is exact.

\paragraph{Contributions.}
\begin{enumerate}[leftmargin=1.4em,itemsep=2pt,topsep=2pt]
  \item \textbf{A pipeline that connects concept definition to certified
  reachability.} On cached GPT-2 (residual layer~6) we define an eval-aware target
  set $A^\star$ from an NLA-style verbalizer and a text classifier, distil it into
  a certifiable ReLU surrogate, and certify reachability of $A^\star$ around benign
  activations. The whole loop runs on a laptop and ships with tests.
  \item \textbf{A sound, self-contained certificate.} We found the published
  \texttt{auto\_LiRPA} library~\cite{xu2020autolirpa} unusable under current
  \texttt{torch}, so we re-derive a CROWN backward bound~\cite{zhang2018crown}
  combined with interval propagation~\cite{gowal2018ibp} from scratch in pure
  \texttt{torch}. Soundness is enforced at runtime: every bound is checked against
  a brute-force sampling witness, and a violation raises immediately.
  \item \textbf{The certificate as a training signal.} Because the bound is
  differentiable, descending it pushes $A^\star$ away from benign activations and
  enlarges the certified radius --- Raghunathan's original move --- which we show as
  a before/after mechanism demonstration.
  \item \textbf{Constructive inversion plus a loop-closure check.} A GCG/HotFlip
  token search~\cite{zou2023gcg,ebrahimi2018hotflip} produces a real prompt whose
  activation crosses the threshold, and a re-verbalization step checks that the
  crossing is semantically real and not a numerical artifact.
  \item \textbf{A robust negative result.} Budget-projected ascent on the surrogate
  can satisfy the threshold, but only far off the activation manifold, where the
  surrogate says ``eval'' and the language readout says ``ordinary.'' This
  off-manifold certificate-gaming is exactly what the loop-closure check is for, and
  it reproduces across anchors.
\end{enumerate}

We are deliberately explicit throughout about the boundary between what is real and
sound (the GPT-2 activations, the certificate, the inversion) and what is an honest
stand-in (the synthetic eval-vs-deploy labels and the retrieval-based verbalizer).
We present the system as a validated methods skeleton rather than a finished
empirical result about real eval-awareness, and we flag four extensions, currently
in progress, that are needed to make it honest science (Section~\ref{sec:results}).
